\section{Kết luận}
% ================= FRAME 1: THÁCH THỨC & GIẢI PHÁP =================
\begin{frame}{Các vấn đề thách thức và Giải pháp}
    \begin{itemize}
        \item<1-> \textbf{1. Hiệu ứng chân trời (Horizon Effect):}
        \begin{itemize}
            \item \textit{Vấn đề:} Engine thí quân vô nghĩa để đẩy đòn bất lợi ra khỏi "chân trời" tìm kiếm, dẫn đến đánh giá sai lầm.
            \item \textit{Giải pháp:} Triển khai \textbf{Quiescence Search} kết hợp Check Extensions. Tiếp tục duyệt các nước ăn quân đến khi thế trận tĩnh hoàn toàn.
        \end{itemize}
        
        \item<2-> \textbf{2. Tranh chấp đồng bộ hóa (UCI):}
        \begin{itemize}
            \item \textit{Vấn đề:} Lệnh ngắt \texttt{stop} từ GUI gửi đến bất kỳ lúc nào. Khóa Mutex truyền thống gây nghẽn hoặc treo luồng.
            \item \textit{Giải pháp:} Dùng biến nguyên tử \textbf{\texttt{AtomicBool}} chia sẻ giữa luồng đọc và luồng tìm kiếm $\rightarrow$ Dừng tức thời ($<1$ ms) không chặn luồng.
        \end{itemize}
    \end{itemize}
\end{frame}

% ================= FRAME 2: CÔNG NGHỆ SỬ DỤNG =================
\begin{frame}{Công nghệ và Tài nguyên sử dụng}
    \textbf{Dự án được xây dựng trên nền tảng công nghệ mã nguồn mở:}
    \vspace{0.2cm}
    \begin{itemize}
        \item<1-> \textbf{Ngôn ngữ \& Công cụ:} \textbf{Rust} (\texttt{rustc}, \texttt{cargo}). Biên dịch tối ưu với \texttt{opt-level=3} và \texttt{target-cpu=native}.
        \item<2-> \textbf{Quản lý bộ nhớ:} Thư viện \textbf{\texttt{arrayvec}} lưu danh sách nước đi trực tiếp trên bộ nhớ Stack, loại bỏ overhead cấp phát Heap động.
        \item<3-> \textbf{Hệ sinh thái Kiểm thử:} 
        \begin{itemize}
            \item \textbf{Python Suite:} Viết kịch bản tự động hóa giải đấu.
            \item \textbf{Sylvan CLI \& GUI:} Giao diện điều phối qua giao thức UCI.
            \item \textbf{Đối chiếu \& Benchmark:} \textit{Fairy-Stockfish}, \textit{Pikafish} (đo Elo) và kiến trúc \textit{Reckless}, \textit{Viridithas} (tham khảo cấu trúc).
        \end{itemize}
    \end{itemize}
\end{frame}

% ================= FRAME 4: SOFTMAX SAMPLING =================
\begin{frame}{Hướng phát triển 1: Softmax Sampling}
    \begin{itemize}
        \item \textbf{Vấn đề:} Việc luôn chọn nước đi tốt nhất (\texttt{bestmove}) khiến lối chơi rập khuôn, dễ bị đối thủ "bắt bài".
        \item \textbf{Giải pháp:} Chuyển điểm số thành phân phối xác suất dựa trên nhiệt độ $\tau$ để đa dạng hóa nước đi:
        $$P(a_i | s) = \frac{e^{Q(s, a_i) / \tau}}{\sum_j e^{Q(s, a_j) / \tau}}$$
    \end{itemize}
    
    \vspace{-1cm}
    \begin{figure}[htbp]
        \centering
        \begin{tikzpicture}[scale=0.65, transform shape]
            % Axes
            \draw[thick, ->] (0,0) -- (8,0) node[right] {Nước đi (Chất lượng giảm dần)};
            \draw[thick, ->] (0,0) -- (0,5) node[above] {Xác suất chọn $P(a_i)$};
            
            % Argmax (tau -> 0)
            \draw[blue, ultra thick] (1, 4.5) -- (1, 0);
            \fill[blue] (1, 4.5) circle (2.5pt);
            \node[blue, right] at (1.1, 4.5) {$\tau \to 0$ (Argmax)};
            
            % Softmax moderate (tau = 0.5)
            \draw[red, thick, smooth] plot coordinates {(0.5, 0.5) (1, 3.5) (2, 1.5) (3, 0.8) (4, 0.4) (6, 0.1)};
            \node[red, right] at (2.1, 1.8) {$\tau = 0.5$ (Cân bằng)};
            
            % Softmax high (tau -> inf)
            \draw[green!60!black, thick, dashed, smooth] plot coordinates {(0.5, 1.2) (1, 1.1) (3, 0.9) (5, 0.8) (7, 0.7)};
            \node[green!60!black, right] at (5.1, 1.0) {$\tau \to \infty$ (Ngẫu nhiên)};
            
            % Labels on X-axis
            \draw (1, 0.1) -- (1, -0.1) node[below] {Tốt nhất};
            \draw (3, 0.1) -- (3, -0.1) node[below] {Tốt nhì};
            \draw (5, 0.1) -- (5, -0.1) node[below] {Bình thường};
        \end{tikzpicture}
    \end{figure}
\end{frame}

% ================= FRAME 5: MẠNG NƠ-RON NNUE =================
\begin{frame}{Hướng phát triển 2: Tích hợp mạng NNUE}
    \begin{columns}[T]
        \begin{column}{0.45\textwidth}
            \vspace{0.2cm}
            \begin{itemize}
                \item Mạng NNUE (Efficiently Updatable Neural Networks) giúp nắm bắt tương tác phi tuyến (ví dụ: Xe-Pháo-Mã).
                \item \textbf{Điểm mạnh:} Cập nhật vi phân. Chỉ tính lại trọng số cho quân cờ vừa di chuyển, giữ vững tốc độ duyệt NPS.
            \end{itemize}
        \end{column}
        
        \begin{column}{0.55\textwidth}
            \centering
            \begin{figure}[htbp]
                \begin{tikzpicture}[
                    scale=0.6, transform shape,
                    node distance=1.5cm and 2cm,
                    neuron/.style={circle, draw, minimum size=0.8cm, inner sep=0pt},
                    active/.style={fill=red!30},
                    hidden/.style={fill=blue!20},
                    output/.style={fill=green!30}
                ]
                    % Input Layer
                    \node[neuron] (i1) at (0, 3) {};
                    \node[neuron, active] (i2) at (0, 2) {};
                    \node[neuron] (i3) at (0, 1) {};
                    \node at (0, 0.3) {$\vdots$};
                    \node[neuron, active] (in) at (0, -1) {};
                    \node[above=0.5cm of i1, text width=2.5cm, align=center] {Lớp đầu vào\\(Đặc trưng thưa)};
                    \node[left=0.2cm of i2] {Quân đi $\to$};
                    
                    % Accumulator
                    \node[neuron, hidden] (a1) at (4, 2.5) {};
                    \node[neuron, hidden] (a2) at (4, 1.5) {};
                    \node[neuron, hidden] (a3) at (4, 0.5) {};
                    \node at (4, -0.2) {$\vdots$};
                    \node[neuron, hidden] (am) at (4, -1) {};
                    \node[above=1cm of a1, text width=3cm, align=center] {Bộ tích lũy\\(Cập nhật vi phân)};
                    
                    % Output Layer
                    \node[neuron, output] (out) at (8, 0.75) {};
                    \node[above=0.5cm of out, text width=2cm, align=center] {Điểm số};
                    
                    % Connections
                    \foreach \i in {i2, in} {
                        \foreach \a in {a1, a2, a3, am} { \draw[->, thick, red] (\i) -- (\a); }
                    }
                    \draw[->, gray, dashed] (i1) -- (a1);
                    \draw[->, gray, dashed] (i3) -- (a2);
                    \foreach \a in {a1, a2, a3, am} { \draw[->, thick, blue] (\a) -- (out); }
                \end{tikzpicture}
            \end{figure}
        \end{column}
    \end{columns}
\end{frame}

% ================= FRAME 6: THUẬT TOÁN SPSA =================
\begin{frame}{Hướng phát triển 3: Tối ưu hóa SPSA}
    \begin{itemize}
        \item \textbf{Vấn đề:} Hàng trăm tham số cắt tỉa tìm kiếm (LMR, NMP) đang được chỉnh tay, không thể dùng Gradient Descent vì đây là hộp đen.
        \item \textbf{Giải pháp:} Dùng thuật toán \textbf{SPSA}. Tạo nhiễu ngẫu nhiên lên toàn bộ tham số, cho engine tự đấu (Self-play) để ước lượng Gradient.
    \end{itemize}
    
    \vspace{-0.25cm}
    \begin{figure}[htbp]
        \centering
        \begin{tikzpicture}[
            scale=0.6, transform shape,
            node distance=1.5cm and 2cm,
            box/.style={rectangle, draw, rounded corners, minimum height=1.2cm, text centered, text width=3.8cm},
            box_sq/.style={rectangle, draw, minimum height=1.2cm, text centered, text width=3.5cm},
            arrow/.style={thick, ->, >=stealth}
        ]
            \node[box, fill=green!20, text width=4cm] (theta) at (0, 4) {Tham số hiện tại $\theta$};
            \node[box, fill=blue!10] (plus) at (-3, 2) {$\theta + c_k\Delta_k$};
            \node[box, fill=blue!10] (minus) at (3, 2) {$\theta - c_k\Delta_k$};
            \node[box_sq, fill=orange!20] (game) at (0, 0) {Tự đấu\\(Engine Match)};
            \node[box, fill=pink!20, text width=4.5cm] (grad) at (0, -2) {Ước lượng\\Gradient $\hat{g}_k(\theta)$};
            
            \draw[arrow] (theta.south) -- (plus.north) node[midway, left, yshift=0.15cm] {+ Nhiễu $\Delta_k$};
            \draw[arrow] (theta.south) -- (minus.north) node[midway, right, yshift=0.15cm] {- Nhiễu $\Delta_k$};
            \draw[arrow] (plus.south) -- (game.north west);
            \draw[arrow] (minus.south) -- (game.north east);
            \draw[arrow] (game.south) -- (grad.north) node[midway, right=0.1cm] {Thắng/Thua/Hòa};
            
            % Feedback loop
            \draw[arrow, rounded corners] (grad.west) -- (-5.5, -2) -- (-5.5, 4) -- (theta.west);
            \node[anchor=south west, align=left] at (-5.7, 4.1) {Cập nhật $\theta_{k+1}$\\$\theta_k - a_k \hat{g}_k(\theta)$};
        \end{tikzpicture}
    \end{figure}
\end{frame}

\begin{frame}{~}
  \centering
  \vfill
  {\Huge \textcolor{hustred}{Cảm ơn mọi người đã lắng nghe!}}

  \vspace{1.2em}
  {\large Q \& A}
  \vfill
\end{frame}


